% ══════════════════════════════════════════════════════════════════════
%  Beamer Presentation
%  On How Voting Mechanisms Affect Welfare and Equity
%  for Participatory Budgeting: An Approach from Durango, Mexico
%
%  Audience:  Policy / mixed (municipal government, civil society)
%  Length:    ~45 minutes (~35 slides)
%  Compile:   latexmk -pdf main.tex  (from presentation/)
% ══════════════════════════════════════════════════════════════════════
\documentclass[12pt, aspectratio=169]{beamer}

% ── Theme ────────────────────────────────────────────────────────────
\usetheme[progressbar=frametitle, block=fill]{metropolis}
\usepackage{appendixnumberbeamer}

% ── Colors ───────────────────────────────────────────────────────────
\definecolor{ujedblue}{HTML}{003580}
\definecolor{ujeddark}{HTML}{001F4D}
\definecolor{ujedaccent}{HTML}{D4A017}
\definecolor{mesdark}{HTML}{1A6B3A}
\definecolor{demdark}{HTML}{8B1A1A}
\definecolor{uamdark}{HTML}{4A4A8A}

\setbeamercolor{palette primary}{bg=ujedblue, fg=white}
\setbeamercolor{progress bar}{fg=ujedaccent, bg=ujeddark}
\setbeamercolor{title separator}{fg=ujedaccent}
\setbeamercolor{frametitle}{bg=ujedblue, fg=white}
\setbeamercolor{alerted text}{fg=demdark}
\setbeamercolor{example text}{fg=mesdark}

% DEM/UAM/MES color helpers
\newcommand{\cDEM}[1]{{\color{demdark}#1}}
\newcommand{\cUAM}[1]{{\color{uamdark}#1}}
\newcommand{\cMES}[1]{{\color{mesdark}#1}}

% ── Packages ─────────────────────────────────────────────────────────
\usepackage{amsmath,amssymb,mathtools}
\usepackage[T1]{fontenc}
\usepackage[utf8]{inputenc}
\usepackage{booktabs}
\usepackage{multirow}
\usepackage{tikz}
\usetikzlibrary{arrows.meta,positioning,calc,decorations.pathreplacing}
\usepackage[authoryear,round]{natbib}
\usepackage{pifont}   % for \cmark / \xmark
\newcommand{\cmark}{\ding{51}}
\newcommand{\xmark}{\ding{55}}

% ── Math macros (identical to shadow_paper.tex) ───────────────────────
\newcommand{\voters}{N}
\newcommand{\districts}{\mathcal{D}}
\newcommand{\projects}{\mathcal{P}}
\newcommand{\agents}{N}
\newcommand{\D}{\mathcal{D}}
\newcommand{\Pset}{\mathcal{P}}
\newcommand{\votes}{v}
\newcommand{\budget}{B}
\newcommand{\util}{u}
\newcommand{\info}{\sigma}
\newcommand{\typespace}{\Theta}
\newcommand{\outcomes}{X}
\newcommand{\DEM}{\text{DEM}}
\newcommand{\UAM}{\text{UAM}}
\newcommand{\MES}{\text{MES}}
\newcommand{\R}{\mathbb{R}}
\newcommand{\E}{\mathbb{E}}
\newcommand{\cost}{c}
\newcommand{\vtotal}{V}
\newcommand{\bshare}{b}
\newcommand{\mesrho}{\varrho}
\newcommand{\salience}{\alpha}
\newcommand{\marg}{\mu}
\newcommand{\Iid}{I_i^d}
\newcommand{\Iidbar}{\bar{I}_i^d}
\newcommand{\bdist}{\beta}
\newcommand{\Ap}{A_i}
\newcommand{\Np}{N(p)}
\DeclareMathOperator*{\argmin}{arg\,min}
\DeclareMathOperator*{\argmax}{arg\,max}

% ── Title ─────────────────────────────────────────────────────────────
\title{Voting Mechanisms, Welfare, and Equity\\in Participatory Budgeting}
\subtitle{\textit{Presupuesto Participativo}: An Approach from Durango, Mexico}
\author{Mario}
\institute{FECA-UJED \textbar{} Universidad Juárez del Estado de Durango}
\date{2026}

% ══════════════════════════════════════════════════════════════════════
\begin{document}
% ══════════════════════════════════════════════════════════════════════

\begin{frame}[plain,shrink=20]
  \titlepage
\end{frame}

% ── Outline ──────────────────────────────────────────────────────────
\begin{frame}[shrink=20]{Roadmap}
  \tableofcontents
\end{frame}

% ══════════════════════════════════════════════════════════════════════
\section{Motivation}
% ══════════════════════════════════════════════════════════════════════

\begin{frame}{Democracy Beyond the Ballot Box}
  \begin{columns}[T]
    \column{0.52\textwidth}
      \textbf{Participatory Budgeting}\\
      \textit{(Presupuesto Participativo)}
      \begin{itemize}
        \item Citizens vote on \emph{how} to spend part of the public budget
        \item Born: Porto Alegre, Brazil (1989)
        \item Now in 11\,000+ cities worldwide
        \item Mexico: adopted by many municipalities under \textit{Leyes de Participación Ciudadana}
      \end{itemize}
    \column{0.45\textwidth}
      \textbf{Why does it matter?}
      \begin{itemize}
        \item Citizens gain a direct say in fiscal decisions
        \item Aligns government incentives with community needs
        \item Good allocation $\to$ political goodwill $\to$ accountability
      \end{itemize}
      \bigskip
      \begin{alertblock}{}
        \centering\small
        \textbf{Key insight:} The outcomes depend critically on the \emph{voting rules}.
      \end{alertblock}
  \end{columns}
\end{frame}

\begin{frame}{Durango's \textit{Presupuesto Participativo}: The Current Rules}
  \begin{block}{The District-Exclusivity Rule (\textit{Regla de Exclusividad Distrital})}
    Each voter must select \textbf{one district} (\textit{sector}) and allocate \textbf{all} votes within it.
    Voting across district boundaries is \textbf{prohibited}.
  \end{block}
  \medskip
  \begin{columns}[T]
    \column{0.48\textwidth}
      \textbf{Mechanism in practice:}
      \begin{itemize}
        \item $n$ voters, each with $\votes = 6$ votes
        \item $m$ districts in Durango
        \item Projects registered within each district
        \item Top vote-getters funded greedily until budget exhausted
      \end{itemize}
    \column{0.48\textwidth}
      \textbf{Stated justification:}
      \begin{itemize}
        \item Prevent high-visibility city-center projects from absorbing all votes
        \item Protect marginalized \textit{colonias} from being shut out
        \item Ensure geographic representation
      \end{itemize}
  \end{columns}
\end{frame}

\begin{frame}{The Central Question}
  \begin{center}
    \Large
    Does the district-exclusivity rule achieve its goals?\\[0.8em]
    \normalsize
    And if not --- what should we use instead?
  \end{center}
  \bigskip
  \textbf{This paper:}
  \begin{enumerate}
    \item Formally identifies \textbf{five distortions} caused by the current rule
    \item Shows that simply removing the constraint \textbf{destroys equity}
    \item Proposes the \textbf{Method of Equal Shares} (\textit{Método de Partes Iguales}) as a reform
    \item Characterizes when and why MES achieves both efficiency \emph{and} equity
  \end{enumerate}
\end{frame}

% ══════════════════════════════════════════════════════════════════════
\section{A Brief Framework: Mechanism Design}
% ══════════════════════════════════════════════════════════════════════

\begin{frame}{What Is Mechanism Design?}
  \begin{quote}
    ``Engineering the rules of the game so that self-interested agents produce good social outcomes.''
  \end{quote}
  \bigskip
  \begin{columns}[T]
    \column{0.32\textwidth}
      \centering\textbf{(1) Environment}\\[0.4em]
      Who are the players?\\
      What do they want?\\
      What are the constraints?
    \column{0.32\textwidth}
      \centering\textbf{(2) Social Choice Function}\\[0.4em]
      What outcome \emph{should} we produce given everyone's preferences?
    \column{0.32\textwidth}
      \centering\textbf{(3) Mechanism}\\[0.4em]
      What \emph{rules} (message space + outcome function) produce that outcome?
  \end{columns}
  \bigskip
  \begin{alertblock}{The design problem}
    Preferences are \emph{private}. How do we design rules so that agents reveal them truthfully?
  \end{alertblock}
\end{frame}

\begin{frame}{Key Definitions}
  \begin{block}{Environment $\langle \agents, \outcomes, \typespace, u \rangle$}
    $n$ agents with private types $\theta_i \in \typespace_i$ (preferences + information);
    utility functions $u_i : \outcomes \times \typespace_i \to \mathbb{R}$.
  \end{block}
  \begin{block}{Social Choice Function $f : \typespace \to \outcomes$}
    The planner's ideal rule: maps type profiles to outcomes.
    \emph{Example:} fund the set of projects with highest total welfare.
  \end{block}
  \begin{block}{Mechanism $(\mathcal{M}, g)$}
    $\mathcal{M} = \prod_i M_i$ is the \textbf{message space} (what agents report);
    $g : \mathcal{M} \to \outcomes$ is the \textbf{outcome function} (how messages map to results).
  \end{block}
\end{frame}

\begin{frame}{The Revelation Principle}
  \begin{theorem}[Revelation Principle]
    If a social choice function $f$ can be \emph{implemented} by \emph{any} mechanism in Nash equilibrium,
    it can also be implemented by a \textbf{direct, incentive-compatible} mechanism ---
    one where truthful reporting is a dominant strategy.
  \end{theorem}
  \bigskip
  \textbf{Implication for us:}
  We only need to study \emph{direct} mechanisms (voters report preferences directly).
  Restricting to truthful reporting is without loss of generality.
\end{frame}

\begin{frame}[shrink=15]{The Gibbard-Satterthwaite Impossibility}
  \begin{theorem}[Gibbard 1973; Satterthwaite 1975]
    With three or more alternatives and \emph{unrestricted} preferences, any
    \textbf{strategy-proof} social choice function that can select any alternative
    must be \textbf{dictatorial}.
  \end{theorem}
  \bigskip
  \begin{itemize}
    \item No non-dictatorial mechanism is fully strategy-proof in a general setting
    \item The \textbf{goal shifts}: design mechanisms that are \emph{as good as possible} under this constraint
    \item We will revisit this when we evaluate the Method of Equal Shares
  \end{itemize}
  \medskip
  \begin{alertblock}{}
    Both Durango's current rule \emph{and} the proposed alternative fail strategy-proofness.\\
    The difference lies in everything else.
  \end{alertblock}
\end{frame}

% ══════════════════════════════════════════════════════════════════════
\section{The Model}
% ══════════════════════════════════════════════════════════════════════

\begin{frame}[shrink=10]{The Participatory Budgeting Environment}
  \begin{block}{Definition: PB Environment $\langle \agents, \D, \Pset, \budget, \votes, \util, \info \rangle$}
  \end{block}
  \begin{columns}[T]
    \column{0.50\textwidth}
      \begin{itemize}
        \item $\agents = \{1,\ldots,n\}$: voters
        \item $\D = \{1,\ldots,m\}$: districts (\textit{sectores})
        \item $\Pset = \bigsqcup_{d} \Pset^d$: projects, partitioned by district
        \item $\cost : \Pset \to \R_{>0}$: cost function
        \item $\budget > 0$: total budget
      \end{itemize}
    \column{0.48\textwidth}
      \begin{itemize}
        \item $\votes \in \mathbb{N}$: vote endowment per voter
        \item $\util_i(p) \geq 0$: voter $i$'s utility for project $p$
        \item \alert{$\info_i(p) \in \{0,1\}$}: \textbf{information function}\\
          ``Does voter $i$ know about project $p$?''
      \end{itemize}
  \end{columns}
  \medskip
  \begin{exampleblock}{Why $\info_i$?}
    Voters know their neighborhood's projects --- not the whole city's.
    This \textbf{rational ignorance} \citep{downs1957} is the root cause of several distortions.
  \end{exampleblock}
\end{frame}

\begin{frame}{District Salience \textit{(Visibilidad Distrital)}}
  \begin{block}{Definition: Salience $\salience_d$}
    Fraction of voters informed about \emph{at least one} project in district $d$:
    \[
      \salience_d = \frac{|\{i \in \agents : \exists\, p \in \Pset^d,\; \info_i(p) = 1\}|}{n}
    \]
  \end{block}
  \bigskip
  \begin{columns}[T]
    \column{0.48\textwidth}
      \textbf{High-salience} $(\salience_d \approx 1)$:\\
      City center, university district.\\
      Most voters know the projects.
    \column{0.48\textwidth}
      \textbf{Low-salience} $(\salience_d \approx 0)$:\\
      Peripheral, marginalized \textit{colonias}.\\
      Few voters know the projects.
  \end{columns}
  \bigskip
  Salience drives both \textbf{strategic behavior} (who goes where) and \textbf{equity outcomes} (who gets funded).
\end{frame}

\begin{frame}[shrink=10]{Two Mechanisms: DEM and UAM}
  \begin{columns}[T]
    \column{0.48\textwidth}
      \begin{block}{\cDEM{District-Exclusivity Mechanism (DEM)}}
        \begin{itemize}
          \item Voter selects \textbf{one district} $d_i \in \D$
          \item Allocates all $\votes$ votes \emph{within} $\Pset^{d_i}$
          \item \textbf{Durango's current rule}
        \end{itemize}
      \end{block}
    \column{0.48\textwidth}
      \begin{block}{\cUAM{Unrestricted Allocation Mechanism (UAM)}}
        \begin{itemize}
          \item Allocates $\votes$ votes across \textbf{any} project in \textbf{any} district
          \item No geographic constraint
          \item The ``obvious'' alternative
        \end{itemize}
      \end{block}
  \end{columns}
  \bigskip
  Both mechanisms use the same \textbf{greedy outcome rule}: fund projects in decreasing vote order until the budget is exhausted.
  \[
    \mathcal{M}^{\cDEM{\DEM}} \;\subsetneq\; \mathcal{M}^{\cUAM{\UAM}}
  \]
  Every DEM vote profile is a valid UAM profile.
  The DEM simply \emph{restricts} the feasible space.
\end{frame}

% ══════════════════════════════════════════════════════════════════════
\section{Five Distortions of the Current Mechanism}
% ══════════════════════════════════════════════════════════════════════

\begin{frame}[shrink=15]{Distortion 1: Preference Truncation \textit{(Truncamiento de Preferencias)}}
  \begin{block}{Proposition}
    Under sincere voting \emph{with full information}, utilitarian social welfare satisfies
    \(W^{\cDEM{\DEM}} \leq W^{\cUAM{\UAM}}\),
    with strict inequality whenever the marginal voter has cross-district preferences.\\[2pt]
    \footnotesize\emph{(With incomplete information the ranking can reverse --- this is exactly Distortion 5.)}
  \end{block}
  \medskip
  \textbf{Intuition:}
  \begin{itemize}
    \item A voter who cares about projects in \emph{two} districts must abandon half their preferences
    \item All 6 votes are forced into one district --- even if the voter's top priority lies elsewhere
    \item The DEM \emph{artificially narrows} each voter's expressed preference profile
    \item Since $\mathcal{M}^{\DEM} \subsetneq \mathcal{M}^{\UAM}$, the DEM truncates the preferences each voter can express
  \end{itemize}
\end{frame}

\begin{frame}[shrink=10]{Distortion 2: Noise Injection \textit{(Inyección de Ruido)}}
  \begin{block}{Noise Allocation Model}
    A voter selecting district $d$ must allocate \emph{all} $\votes$ votes within $\Pset^d$.
    Votes on projects the voter doesn't know become \textbf{noise}: random allocation unrelated to quality.
  \end{block}
  \begin{block}{Proposition}
    For any project $p \in \Pset^d$:
    \[
      \E[\vtotal(p)] \;=\; \underbrace{S(p)}_{\text{signal: votes from informed voters}} + \underbrace{\E[\eta(p)]}_{\text{noise: votes from uninformed voters}}
    \]
    $\E[\eta(p)]$ is \textbf{independent of $\util_j(p)$} for all voters $j$.
    The noise-to-signal ratio rises with the number of unknown projects.
  \end{block}
\end{frame}

\begin{frame}[shrink=15]{Noise in Action: Durango-Scale Example}
  \textbf{Setup:} District $A$, 8 projects, 100 voters, $\votes = 6$
  \begin{itemize}
    \item \textbf{40 institutional voters}: know only $p_1$ (flagship project)
    \item \textbf{30 local residents}: know 3 projects including $p_1$
    \item \textbf{30 mobilized supporters}: each knows one other project (not $p_1$)
  \end{itemize}
  \medskip
  \begin{columns}
    \column{0.48\textwidth}
      \textbf{Dominant project $p_1$:}\\[0.3em]
      Signal\quad $S(p_1) \approx 220$\\
      Noise\quad $\E[\eta(p_1)] \approx 4.3$\\[0.3em]
      \textcolor{mesdark}{Noise/Signal $\approx \mathbf{2\%}$\quad\cmark}
    \column{0.48\textwidth}
      \textbf{Small valued project $p_7$:}\\[0.3em]
      Signal\quad $S(p_7) = 10$\\
      Noise\quad $\E[\eta(p_7)] \approx 15.7$\\[0.3em]
      \alert{Noise/Signal $\approx \mathbf{157\%}$\quad\xmark}
  \end{columns}
  \bigskip
  \begin{alertblock}{}
    For marginal projects, noise \emph{exceeds} the signal.
    Funding at the margin is driven by random allocation, not genuine preferences.
  \end{alertblock}
\end{frame}

\begin{frame}[shrink=15]{Distortion 3: The Coordination Game \textit{(Juego de Coordinación)}}
  \begin{block}{District-Selection Game $\Gamma$}
    Voters with positive utility in multiple districts face a strategic choice: which district to select?
    Payoffs depend on others' choices --- creating \textbf{strategic complementarity}.
  \end{block}
  \begin{block}{Proposition}
    If enough voters have cross-district preferences, $\Gamma$ has \textbf{at least two Nash equilibria}
    producing \textbf{distinct funded sets} with \textbf{distinct welfare levels}.
  \end{block}
  \medskip
  \textbf{Example:} Districts $A$ and $B$ each have one borderline project; 50 mobile voters can tip either.
  \begin{itemize}
    \item $E_A$: all mobile voters $\to A$ \;---\; $p_A$ funded, $p_B$ not
    \item $E_B$: all mobile voters $\to B$ \;---\; $p_B$ funded, $p_A$ not
  \end{itemize}
  \alert{Which project gets funded depends on \emph{who shouts louder}, not on what the community values.}
\end{frame}

\begin{frame}[shrink=15]{Distortion 4: Mobilization Advantage \textit{(Ventaja por Movilización)}}
  \begin{block}{Proposition}
    A mobilizing institution $(\mathcal{C}, p_\mathcal{C})$ coordinating $|\mathcal{C}|$ voters:
    \begin{enumerate}[\upshape(i)]
      \item Inflates district $d_\mathcal{C}$'s vote pool by $|\mathcal{C}| \cdot \votes$
      \item Creates vote displacement $\Delta_\mathcal{C} = \votes \cdot |\mathcal{C}^*|$
        ($\mathcal{C}^*$: members diverted from their preferred district)
      \item Causes private utility loss to diverted members, $\mathcal{L}(\mathcal{C}) \geq 0$
    \end{enumerate}
  \end{block}
  \medskip
  \textbf{Example:} A university coordinates 400 employees $\to$ the district with the university park project.
  \begin{itemize}
    \item $\Delta_\mathcal{C} = 400 \times 6 = 2{,}400$ displaced votes
    \item Those 400 voters can no longer support the projects they actually value most
    \item The outcome reflects institutional power, not community preferences
  \end{itemize}
\end{frame}

\begin{frame}[shrink=10]{Distortion 5: The Equity Tradeoff \textit{(Dilema de Equidad)}}
  \begin{block}{Proposition}
    \begin{itemize}
      \item \cUAM{\textbf{UAM}}: High-salience districts capture nearly all votes.
        Marginalized, low-salience districts may get \emph{zero funding} ($\beta$-equity fails with $\beta = 0$).
      \item \cDEM{\textbf{DEM}}: Committed local residents create a \emph{protected vote pool},
        giving marginalized districts partial protection ($\beta > 0$).
    \end{itemize}
  \end{block}
  \medskip
  \begin{alertblock}{The core dilemma}
    The DEM's district constraint is a \textbf{blunt} equity instrument.\\
    It partially protects marginalized districts --- but at the cost of distortions 1--4.\\[0.5em]
    \textbf{Simply removing the constraint (UAM) makes equity \emph{worse.}}\\[0.5em]
    We need a mechanism that achieves equity \emph{by design}, not by restriction.
  \end{alertblock}
\end{frame}

% ══════════════════════════════════════════════════════════════════════
\section{The Method of Equal Shares}
% ══════════════════════════════════════════════════════════════════════

\begin{frame}[shrink=15]{The Core Idea: Individual Budget Shares}
  \begin{center}
    \Large
    Each voter receives an \textbf{equal share} of the budget:\\[0.3em]
    $\bshare_i = \budget / n$\\[0.5em]
    \normalsize
    Projects are funded by \textbf{pooling} the shares of their supporters.
  \end{center}
  \medskip
  \begin{columns}[T]
    \column{0.48\textwidth}
      \textbf{What voters do:}
      \begin{itemize}
        \item Submit an \textbf{approval ballot} $\Ap \subseteq \Pset$
        \item Approve any projects they support, \emph{across any district}
        \item Only approve projects they know: $\info_i(p) = 1$
      \end{itemize}
    \column{0.48\textwidth}
      \textbf{What the algorithm does:}
      \begin{itemize}
        \item Finds the project with the \emph{most support relative to cost}
        \item Deducts from supporters' shares proportionally
        \item Repeats until no more projects are affordable
      \end{itemize}
  \end{columns}
\end{frame}

\begin{frame}{The MES Algorithm}
  \begin{block}{Definition: Method of Equal Shares}
    \textbf{Initialize:} $\bshare_i = \budget/n$ for all $i$;\quad $W = \emptyset$.\\[0.5em]
    \textbf{Repeat until no project is affordable:}
    \begin{enumerate}
      \item For each $p \notin W$, find supporters $\Np = \{i : p \in \Ap,\; \bshare_i > 0\}$
      \item Compute the \textbf{cost-per-supporter threshold}:
        \[
          \mesrho(p) \;=\; \min\!\Bigl\{\mesrho \geq 0 :
            \textstyle\sum_{i \in \Np} \min(\bshare_i,\, \mesrho) \geq \cost(p)\Bigr\}
        \]
      \item \textbf{Fund} $p^* = \argmin_p \mesrho(p)$;\quad deduct $\min(\bshare_i, \mesrho(p^*))$ from each $i \in N(p^*)$
    \end{enumerate}
    \textbf{Completion:} Allocate remaining budget greedily by total approval count.
  \end{block}
\end{frame}

\begin{frame}{MES Step-by-Step: Example}
  \textbf{4 voters, $\budget = 100$, $\bshare_i = 25$ each}
  \begin{center}
  \small
  \begin{tabular}{lcccccc}
    \toprule
    Project & Cost & District & $i{=}1$ & $i{=}2$ & $i{=}3$ & $i{=}4$ \\
    \midrule
    $p_1$: Park renovation   & 40 & Center   & \cmark & \cmark & \cmark & \\
    $p_2$: School repair     & 30 & Marginal & & & \cmark & \cmark \\
    $p_3$: Road paving       & 50 & Center   & \cmark & \cmark & & \\
    $p_4$: Water access      & 20 & Marginal & & & & \cmark \\
    \bottomrule
  \end{tabular}
  \end{center}
  \medskip
  \textbf{Round 1} --- thresholds: $\mesrho(p_1)=13.3$, $\mesrho(p_2)=15$,
    $\mesrho(p_3)=25$, $\mesrho(p_4)=20$.\\
  \textcolor{mesdark}{$\to$ Fund $p_1$.} Deduct 13.3 from $i=1,2,3$.
  Remaining: $\bshare_{1,2,3} = 11.67$;\; $\bshare_4 = 25$.
\end{frame}

\begin{frame}[shrink=15]{MES Step-by-Step: Example (cont.)}
  After Round 1: $\bshare_1 = \bshare_2 = \bshare_3 = 11.67$;\quad $\bshare_4 = 25$

  \medskip
  \textbf{Round 2} --- $p_2$: $\mesrho(p_2) = 18.3$;\quad $p_4$: $\mesrho(p_4) = 20$.\\
  $p_3$: total available $= 2 \times 11.67 = 23.3 < 50$ --- \emph{not affordable}.\\
  \textcolor{mesdark}{$\to$ Fund $p_2$ (smallest $\mesrho$).} $i=3$ pays 11.67, $i=4$ pays 18.33.
  Now $\bshare_3 = 0$;\; $\bshare_4 = 6.67$.

  \medskip
  \textbf{Round 3} --- $p_4$: only $i=4$ left, $\bshare_4 = 6.67 < 20$ --- not affordable.\\
  $p_3$: not affordable.\quad \textbf{Stop.}

  \medskip
  \textbf{Completion:} Remaining budget $= 30$.
  Fund $p_4$ (cost 20) greedily. Final: $W = \{p_1, p_2, p_4\}$.

  \medskip
  \begin{exampleblock}{}
    The ``Marginal'' district gets two projects --- without any geographic constraint.\\
    Budget proportionality did the work that the DEM tried to do by force.
  \end{exampleblock}
\end{frame}

\begin{frame}{Extended Justified Representation (EJR)}
  \begin{block}{Definition: EJR}
    A funded set $W$ satisfies \textbf{Extended Justified Representation} if:
    for every cohesive group $S$ of voters whose collective budget share covers their projects,
    at least one member of $S$ has their proportional share of the budget represented in $W$.
  \end{block}
  \begin{theorem}[MES Satisfies EJR \citep{peters2021}]
    The Method of Equal Shares always produces outcomes satisfying EJR.
  \end{theorem}
  \medskip
  \textbf{What this means for Durango:}
  Any community of residents who collectively have enough budget share to fund their projects
  \emph{will} have those interests represented in the outcome --- by construction, not by chance.
\end{frame}

\begin{frame}[shrink=10]{MES Resolves All Four Distortions}
  \begin{columns}[T]
    \column{0.48\textwidth}
      \begin{exampleblock}{Prop.~6: No Preference Truncation}
        Ballot space $M_i^{\MES} = 2^{\Pset}$:
        approve anything, across any district.
        No cross-district restriction.
      \end{exampleblock}
      \begin{exampleblock}{Prop.~7: No Forced Noise}
        Voters approve only projects they know.
        No minimum ballot size.
        Noise votes do not arise.
      \end{exampleblock}
    \column{0.48\textwidth}
      \begin{exampleblock}{Prop.~8: No Coordination Game}
        No district-selection stage.
        The strategic complementarity that creates multiple equilibria is absent by design.
      \end{exampleblock}
      \begin{exampleblock}{Prop.~9: Bounded Mobilization}
        Coalition $\mathcal{C}$ can direct at most $|\mathcal{C}| \cdot (\budget/n)$ to any project.
        Vote displacement $\Delta_\mathcal{C} = 0$.
      \end{exampleblock}
  \end{columns}
\end{frame}

\begin{frame}[shrink=15]{MES and Geographic Equity \textit{(Equidad Geográfica)}}
  \begin{block}{Proposition (Geographic Equity via EJR)}
    Let $\bdist_d = n_d/n$ be district $d$'s population share. If the committed residents of district
    $d$ \emph{commonly approve} a local bundle worth $\bdist_d\budget$, then under MES the funded set
    $W$ satisfies:
    \[
      \sum_{p \in W \cap \Pset^d} \cost(p) \;\geq\; \bdist_d \cdot \budget .
    \]
  \end{block}
  \medskip
  \begin{itemize}
    \item Every such district gets \textbf{at least its proportional budget share} --- guaranteed
    \item This is \textbf{structural} (follows from EJR): holds regardless of what other voters do
    \item Contrast with DEM: equity depends on committed residents happening to form a large enough bloc
    \item If local preferences are heterogeneous, the floor is the largest \emph{commonly-approved}
      bundle --- still $>0$, still beats the UAM's $\beta=0$
  \end{itemize}
\end{frame}

\begin{frame}[shrink=15]{An Honest Limitation: MES Is Not Strategy-Proof}
  \begin{alertblock}{Proposition}
    MES does not satisfy DSIC: there exist preference profiles where a voter gains
    by misreporting their ballot.
  \end{alertblock}
  \medskip
  \begin{itemize}
    \item \textbf{Inherent, not incidental}: \citet{peters2021} (Thm.~4.8) prove that no EJR-satisfying
      mechanism can be strategy-proof
    \item The tradeoff between \emph{proportional representation} and \emph{strategy-proofness}
      cannot be resolved
    \item MES is strategically \emph{simpler} than the DEM: no district selection game,
      no coordination failures, no mobilization arms race
  \end{itemize}
  \medskip
  \begin{block}{}
    The DEM also fails strategy-proofness (its district-selection stage is directly manipulable).\\
    Both mechanisms lose on the same criterion --- MES wins on everything else.
  \end{block}
  \medskip
  \footnotesize\emph{In practice the manipulation needs detailed knowledge of others' ballots and is
  risky --- low-stakes at scale.}
\end{frame}

\begin{frame}[shrink=15]{Two More Honest Limitations \textit{(Otras Limitaciones)}}
  \begin{block}{The completion gap}
    MES often \textbf{stops before spending the whole budget}: once no project is affordable to its
    backers, a \emph{completion} step (greedy by approvals) places the rest --- with \textbf{none} of
    the EJR guarantee.
  \end{block}
  \begin{itemize}
    \item The gap depends on project costs vs.\ the per-voter share $\bshare_i=\budget/n$:
      many small projects $\to$ the proportional phase reaches far; few expensive projects $\to$
      completion dominates
    \item \textbf{Design lever:} keep project cost tiers small relative to $\budget/n$
  \end{itemize}
  \medskip
  \begin{block}{Ballot complexity}
    MES asks voters to consider \emph{all} projects city-wide, not just one district's.
  \end{block}
  \begin{itemize}
    \item With 50+ projects this is a real cognitive load; thin ballots waste a voter's share
    \item \textbf{Mitigation:} a grouped ballot interface with a live ``budget impact'' indicator ---
      richer than the DEM, still tractable
  \end{itemize}
\end{frame}

% ══════════════════════════════════════════════════════════════════════
\section{Comparing the Three Mechanisms}
% ══════════════════════════════════════════════════════════════════════

\begin{frame}[shrink=20]{The Property Table}
  \begin{center}
  \renewcommand{\arraystretch}{1.45}
  \begin{tabular}{p{5.2cm}ccc}
    \toprule
    \textbf{Property} & \cDEM{\textbf{DEM}} & \cUAM{\textbf{UAM}} & \cMES{\textbf{MES}} \\
    \midrule
    No preference truncation    & \xmark & \cmark & \cmark \\
    No forced noise injection   & \xmark & \cmark & \cmark \\
    No coordination game        & \xmark & \cmark & \cmark \\
    Bounded mobilization power  & \xmark & Partial & \cmark \\
    $\beta$-geographic equity   & Partial & \xmark & \cmark \\
    Strategy-proof (DSIC)       & \xmark & \xmark & \xmark \\
    Proportional representation & \xmark & \xmark & \cmark \\
    \bottomrule
  \end{tabular}
  \end{center}
  \medskip
  \cMES{MES \textbf{strictly dominates}} both alternatives on every property
  except strategy-proofness --- which no proportional (EJR) mechanism can satisfy \citep{peters2021}.
\end{frame}

\begin{frame}{300-Voter Comparison: Setup \textit{(Escenario a Escala Durango)}}
  $n = 300$ voters,\quad $\budget = 300$,\quad $\votes = 6$\\[0.5em]
  \begin{center}
  \small
  \begin{tabular}{lccl}
    \toprule
    & Committed residents & Salience $\salience_d$ & Projects \\
    \midrule
    District $A$ (Center)      & 100 & 0.90 & $p_1$ (cost 120),\; $p_2$ (cost 80) \\
    District $B$ (Residential) & 120 & 0.50 & $p_3$ (cost 100) \\
    District $C$ (Marginalized)&  80 & 0.15 & $p_4$ (cost 60),\; $p_5$ (cost 40) \\
    \bottomrule
  \end{tabular}
  \end{center}
  \medskip
  Additional types: 80 mobile voters in $B$ also approve $p_1$;\; 20 connected voters
  in $C$ also approve $p_1$.
\end{frame}

\begin{frame}[shrink=20]{300-Voter Comparison: Results}
  \begin{center}
  \renewcommand{\arraystretch}{1.45}
  \begin{tabular}{p{4.4cm}ccc}
    \toprule
    & \cDEM{\textbf{DEM}} & \cUAM{\textbf{UAM}} & \cMES{\textbf{MES}} \\
    \midrule
    District $A$ funded & $p_1$ & $p_1, p_2$ & $p_1$ \\
    District $B$ funded & $p_3$ & $p_3$ & $p_3$ \\
    District $C$ funded & $p_4$ & \alert{nothing} & $p_5$ \\
    Noise votes present & Yes & No & No \\
    Coordination required & Yes & No & No \\
    Vote displacement & 70 voters & 0 & 0 \\
    \bottomrule
  \end{tabular}
  \end{center}
  \medskip
  \begin{itemize}
    \item \cUAM{UAM}: excludes the marginalized district entirely despite genuine demand
    \item \cDEM{DEM}: partially protects district $C$ but wastes 70 voters' preferences
    \item \cMES{MES}: funds district $C$ through proportionality alone --- \textbf{no constraint needed}
  \end{itemize}
\end{frame}

% ══════════════════════════════════════════════════════════════════════
\section{Conclusion and Policy Implications}
% ══════════════════════════════════════════════════════════════════════

\begin{frame}[shrink=10]{Policy Recommendations \textit{(Recomendaciones de Política)}}
  \begin{enumerate}
    \item \textbf{Replace the DEM with MES} in Durango's \textit{presupuesto participativo}
      \begin{itemize}
        \item Allow cross-district approval voting
        \item Geographic equity is achieved structurally, not by restriction
        \item Implementation is straightforward: voters mark projects they support
      \end{itemize}
    \item \textbf{Invest in project-level information campaigns}
      \begin{itemize}
        \item MES rewards informed voters: more $\info_i(p)=1$ entries $\to$ better outcomes
        \item Targeted publicity for projects in low-salience, marginalized districts
      \end{itemize}
  \end{enumerate}
\end{frame}

\begin{frame}[shrink=15]{Counterfactual Simulation: Durango 2025}
  \textbf{Setup:} 48 projects, 6 districts, \$500{,}000 MXN budget.
  Synthetic approval ballots ($n=1{,}000$, seed 42) generated from observed 2025 vote totals.
  \medskip
  \begin{center}
  \renewcommand{\arraystretch}{1.3}
  \small
  \begin{tabular}{lrr}
    \toprule
    \textbf{Metric} & \cDEM{\textbf{DEM (observed)}} & \cMES{\textbf{MES (counterfactual)}} \\
    \midrule
    Welfare (total votes on funded projects) & 4{,}985 & 4{,}888 \\
    Gini (district-level equity)             & 0.300  & 0.367  \\
    Projects funded                          & 10     & 10     \\
    \midrule
    District 1 funded (MXN) & 50{,}000 & 100{,}000 \\
    District 2 funded (MXN) & 150{,}000 & 150{,}000 \\
    District 3 funded (MXN) & 100{,}000 & 50{,}000 \\
    District 4 funded (MXN) & 0 & 0 \\
    District 5 funded (MXN) & 100{,}000 & 50{,}000 \\
    District 6 funded (MXN) & 100{,}000 & 150{,}000 \\
    \bottomrule
  \end{tabular}
  \end{center}
  \medskip
  \begin{alertblock}{Key finding}
    MES applied to \emph{DEM-constrained} ballots does not outperform DEM.
    Since synthetic voters each select only one district, MES never observes cross-district preferences ---
    the same information gap that drives Propositions~1 and~2.
    The simulation \textbf{illustrates the distortion}, not MES's full potential.
  \end{alertblock}
\end{frame}

\begin{frame}[shrink=20]{Summary \textit{(Resumen)}}
  \begin{enumerate}
    \item The \cDEM{District-Exclusivity Mechanism} introduces \textbf{five distortions}:
      preference truncation, noise injection, coordination failures, mobilization advantage,
      and an equity tradeoff
    \item The ``obvious'' fix (\cUAM{UAM}) \textbf{destroys geographic equity}
      for marginalized districts
    \item The \cMES{Method of Equal Shares} resolves all five distortions while providing
      \textbf{structural} geographic equity via EJR
    \item A counterfactual simulation on Durango 2025 data shows that MES applied to
      DEM-constrained ballots \textbf{replicates the distortion} --- confirming that
      full welfare gains require reforming \emph{both} the mechanism and the ballot format
  \end{enumerate}
  \bigskip
  \begin{center}
    \large\itshape
    Democracy is not only cast in the ballots.\\
    \normalsize
    It is also in the rules that translate ballots into outcomes.
  \end{center}
\end{frame}

\begin{frame}[standout]
  \begin{center}
    Thank you\\[0.8em]
    {\normalsize Gracias}\\[1.5em]
    {\small\ttfamily hola@marionomics.com}
  \end{center}
\end{frame}

% ══════════════════════════════════════════════════════════════════════
\appendix
% ══════════════════════════════════════════════════════════════════════

\begin{frame}[allowframebreaks]{References}
  \bibliographystyle{plainnat}
  \bibliography{references}
\end{frame}

\end{document}
